\chapter{Architecture de la bibliothèque de pricing}\label{ch:pl:pricing-library-architecture}

Un client a demandé vendredi un prix pour un autocallable worst-of à coupons mémoire, avec un lookback sur la barrière de knock-in~: le capital est exposé
dès que la pire performance a été inférieure à 60~\% à une date d'observation, et non seulement à l'échéance. La bibliothèque ne contenait aucune classe
pour cela. Dès lundi, le produit était pricé à partir de onze lignes de \omterm{def:pl:pricing-library-architecture:script}{script de payoff}, sur les mêmes \omterm{def:pl:pricing-library-architecture:market}{objets de marché}, le même modèle et le même moteur
Monte Carlo que toutes les autres transactions de la bibliothèque, avec les grecques, les scénarios et les \omterm{def:pl:the-platform-map:batch}{traitements par lots} en prime. Ce chapitre porte sur l'architecture qui rend cela possible~: une bibliothèque en couches, dont les moteurs ne savent pas quels produits ils
pricent, et un petit langage dans lequel un produit est une donnée.

\section{La bibliothèque analytique et ses couches}

\begin{definition}[Bibliothèque analytique]\label{def:pl:pricing-library-architecture:library}
Une \emph{bibliothèque analytique}\index{bibliothèque analytique} est le code commun de la firme pour valoriser et gérer le risque des instruments~: \omterm{def:pl:pricing-library-architecture:market}{objets de marché},
modèles, moteurs numériques et instruments, séparés de sorte que chacun puisse être remplacé ou étendu sans les autres, et appelés de la même façon par
les outils de pricing, les systèmes de risque et la recherche.
\end{definition}

Le livre~5 en a construit une (\texttt{firm.pricing}, chapitre~28) et en a figé l'interface pour le reste de la série~: les instruments sont des descriptions de
flux de trésorerie~; les données de marché sont un instantané immuable~; un modèle décrit l'évolution du marché~; un moteur calcule la valeur d'un instrument selon
un modèle sur un instantané~; et chaque grecque, chaque scénario et chaque \omterm{def:pl:the-platform-map:batch}{traitement par lots} est une seule fonction de repricing appliquée à des instantanés modifiés. Les couches
font l'architecture~: un nouveau moteur price les anciens instruments, un nouvel instrument utilise les anciens moteurs, et un système de risque (chapitre~18, et le moteur du livre~6)
n'a besoin de connaître que la fonction de repricing.

\begin{omfigure}
\begin{tikzpicture}[font=\scriptsize,box/.style={draw,rounded corners=2pt,align=center,minimum height=0.8cm,inner sep=3pt},>=Stealth]
\node[box,fill=gray!10,minimum width=11cm] (api) at (0,2.2) {appels~: \texttt{price}, \texttt{greeks}, \texttt{sensitivities},
\texttt{reprice}, \texttt{price\_batch} --- une seule fonction de repricing sur des instantanés modifiés};
\node[box,fill=omMeth!12,minimum width=2.6cm] (i) at (-4.7,0.9) {instruments\\ flux de trésorerie};
\node[box,fill=omMeth!12,minimum width=2.6cm] (m) at (-1.57,0.9) {objets de marché\\ courbes, surfaces};
\node[box,fill=omMeth!12,minimum width=2.6cm] (mo) at (1.57,0.9) {modèles\\ Black--Scholes, Heston};
\node[box,fill=omMeth!12,minimum width=2.6cm] (e) at (4.7,0.9) {moteurs\\ analytique, EDP, MC};
\node[box,fill=omDef!15,minimum width=3.8cm] (s) at (-3.2,-0.6) {instrument scripté\\ script + échéanciers};
\node[box,fill=omDef!15,minimum width=3.8cm] (se) at (3.2,-0.6) {moteur de script\\ analyse, échéancier, évaluation};
\draw[->] (s.north) -- (s.north|-i.south); \draw[->] (se.north) -- (se.north|-e.south);
\draw[->] (se.west) -- (s.east);
\foreach \x in {i,m,mo,e} {\draw[->] (\x.north) -- (\x.north|-api.south);}
\node[font=\scriptsize\itshape] at (0,-1.55) {enregistré avec \texttt{firm.pricing.register}~: pricé, choqué et traité par lots comme les autres};
\end{tikzpicture}

\omcaption{Les couches de la bibliothèque et la place du moteur de script dans celles-ci. Un instrument scripté est un instrument~; le moteur de script est un
moteur, enregistré pour le modèle de Black--Scholes~; tout ce qui est au-dessus --- grecques, scénarios, \omterm{def:pl:the-platform-map:batch}{traitements par lots} --- s'applique sans modification.}\label{fig:pl:pricing-library-architecture:layers}
\end{omfigure}

\section{Instruments, objets de marché, modèles et moteurs}

\begin{definition}[Objet de marché]\label{def:pl:pricing-library-architecture:market}
Un \emph{objet de marché}\index{objet de marché} est un objet de la bibliothèque qui représente un élément de l'état du marché --- un cours, une courbe d'actualisation, une
surface de volatilité, une corrélation --- et répond à des questions à son sujet (un facteur d'actualisation, une volatilité implicite), de sorte que les moteurs lisent le
marché à travers une interface fixe, quelle que soit sa construction.
\end{definition}

Dans \texttt{firm.pricing}, un instantané contient les \omterm{def:pl:pricing-library-architecture:market}{objets de marché}, et un choc crée un nouvel instantané dans lequel un objet est remplacé. Le moteur Monte Carlo
lit les forwards, les variances et les corrélations dans cet instantané et génère les trajectoires à partir d'une graine fixée par instrument (nombres aléatoires communs, livre~5), de sorte
qu'un prix choqué ne diffère du prix de base que par le choc seul. Le moteur de script de ce chapitre réutilise ces trajectoires~: c'est le moteur Monte Carlo
dont le payoff est tiré d'un script et non d'une classe.

\section{Scripts de payoff}

\begin{definition}[Langage de script de payoff, script de payoff]\label{def:pl:pricing-library-architecture:script}
Un \emph{langage de script de payoff}\index{langage de script de payoff} est un petit langage de programmation dans lequel les flux de trésorerie d'un produit sont
écrits comme des événements aux dates d'observation --- conditions sur la trajectoire, paiements, état reporté d'une date à l'autre --- et évalués par un
moteur générique sur des trajectoires simulées. Un \emph{script de payoff}\index{script de payoff} est un produit écrit dans ce langage~: une donnée que la bibliothèque price,
et non du code qu'il faudrait étendre.
\end{definition}

Le langage du chapitre (\texttt{firm.payoffdsl}) ne contient que ce dont les produits du chapitre ont besoin~: variables d'état, blocs exécutés aux dates
d'échéanciers nommés, affectations, \texttt{if}/\texttt{else}, \texttt{pay} et \texttt{stop}, arithmétique et comparaisons, ainsi que
\texttt{min}, \texttt{max} et \texttt{abs}. À chaque date, un bloc voit \texttt{S} (le cours au comptant du premier sous-jacent), \texttt{perf} (la pire
performance) et \texttt{t}. Un analyseur lexical transforme le texte en jetons, un analyseur syntaxique par descente récursive les transforme en arbre syntaxique avec une
position d'erreur pour tout élément illisible, et l'échéancier est extrait de l'arbre~: les dates de chaque bloc, dans l'ordre des dates, et dans l'ordre du script
à une même date.

\begin{omfigure}
\begin{tikzpicture}[font=\scriptsize,box/.style={draw,rounded corners=2pt,align=center,minimum height=0.9cm,minimum width=1.9cm,
  fill=omDef!10},>=Stealth]
\node[box] (a) at (0,0) {texte du script};
\node[box] (b) at (2.5,0) {jetons\\ avec positions};
\node[box] (c) at (5.0,0) {arbre\\ syntaxique};
\node[box] (d) at (7.6,0) {échéancier\\ d'événements\\ (date, bloc)};
\node[box,fill=omMeth!12] (e1) at (10.7,0.6) {interpréteur\\ par trajectoire};
\node[box,fill=omMeth!12] (e2) at (10.7,-0.6) {tableaux compilés,\\ masques};
\draw[->] (a) -- (b); \draw[->] (b) -- (c); \draw[->] (c) -- (d); \draw[->] (d) -- (e1.west); \draw[->] (d) -- (e2.west);
\node[font=\scriptsize\itshape,omThm] at (3.4,-0.95) {erreurs avec ligne et colonne};
\node[font=\scriptsize\itshape,align=center] at (7.6,1.15) {trajectoires de la bibliothèque\\ à ces dates};
\end{tikzpicture}

\omcaption{D'un script au prix~: l'analyseur lexical et l'analyseur syntaxique construisent l'arbre, l'échéancier est lu à partir de ses blocs, et l'un ou l'autre
évaluateur fait tourner l'arbre sur les trajectoires de la bibliothèque aux dates de l'échéancier.}\label{fig:pl:pricing-library-architecture:pipeline}
\end{omfigure}

\omcode{code/platforms/19-pricing-library-architecture/python/pl_payoffdsl.py}{36}{44}{L'autocallable Phoenix worst-of à coupons mémoire
sous forme de script de payoff~: coupons à chaque date d'observation, rappel anticipé à toutes les dates sauf la dernière, capital exposé
à l'échéance.}\label{lst:pl:pricing-library-architecture:autocall}

\omcode{code/firm/payoffdsl/firm_payoffdsl.py}{176}{199}{Les instructions de l'analyseur~: \texttt{if}/\texttt{else}, \texttt{pay}, \texttt{stop}
et affectation~; chaque jeton conserve sa ligne et sa colonne pour les messages d'erreur.}\label{lst:pl:pricing-library-architecture:parser}

Le produit du client change en trois lignes~: une variable d'état \texttt{lowest}, mise à jour avec \texttt{lowest = min(lowest, perf)} à chaque
observation, et le test de maturité porte sur \texttt{lowest} au lieu de \texttt{perf}. Il se price à 94.25 pour 100, soit 1.63 de moins que
l'autocallable simple à 95.88~: la barrière de knock-in lookback se déclenche sur les trajectoires qui passent sous 60~\% puis remontent.

\section{Deux façons d'évaluer un script}

L'évaluateur le plus simple parcourt l'arbre syntaxique une fois par trajectoire et par événement, avec des valeurs scalaires~: c'est un interpréteur. Il est facile à vérifier et lent.
L'évaluateur compilé parcourt l'arbre une fois par événement, avec des tableaux sur l'ensemble des trajectoires~: chaque instruction s'exécute sous un masque des trajectoires
auxquelles elle s'applique, un \texttt{if} scinde le masque, un \texttt{pay} ajoute sous ce masque, un \texttt{stop} retire des trajectoires de l'ensemble vivant
(voir le \cref{lst:pl:pricing-library-architecture:compile}). Les deux donnent les mêmes nombres jusqu'au dernier chiffre (les tests les comparent)~; le second
est 19 à 34 fois plus rapide (voir la \cref{fig:pl:pricing-library-architecture:speed}).

\omcode{code/firm/payoffdsl/firm_payoffdsl.py}{357}{388}{L'évaluation compilée~: des masques au lieu de branches, des tableaux au lieu de
boucles.}\label{lst:pl:pricing-library-architecture:compile}

\begin{omfigure}
\begin{tikzpicture}
\begin{loglogaxis}[width=11cm,height=5.6cm,xlabel={trajectoires (échelle log)},ylabel={secondes par prix (échelle log)},grid=major,grid style={gray!20},
  tick label style={font=\small},label style={font=\small},table/col sep=comma,legend style={font=\footnotesize,at={(0.5,1.03)},
  anchor=south,legend columns=2}]
\addplot[color=omThm,very thick,mark=*] table[x=paths,y=interpreted_s]{figdata/platforms/19-pricing-library-architecture/measured_speed.csv};
\addlegendentry{interprété (par trajectoire)}
\addplot[color=omDef,very thick,mark=square*] table[x=paths,y=compiled_s]{figdata/platforms/19-pricing-library-architecture/measured_speed.csv};
\addlegendentry{compilé (tableaux sur les trajectoires)}
\end{loglogaxis}
\end{tikzpicture}

\omcaption{Temps mis pour pricer le script autocallable par l'interpréteur et par l'évaluateur compilé, génération des trajectoires comprise, meilleur de trois exécutions,
un seul thread. La version compilée est 18.8 fois plus rapide à 1\,000 trajectoires, 34.3 fois à 10\,000 et 33.4 fois à 30\,000. Mesures faites sur un portable (Intel Core Ultra 7
155H) sous WSL2, la machine étant par ailleurs inactive. Données~: \texttt{bench\_payoffdsl.py}.}\label{fig:pl:pricing-library-architecture:speed}
\end{omfigure}

Le moteur de script est enregistré dans la bibliothèque comme tout autre moteur (\cref{lst:pl:pricing-library-architecture:engine})~: il analyse le script,
extrait l'échéancier, demande au moteur Monte Carlo des trajectoires aux dates de l'échéancier avec la graine de l'instrument, et évalue. Comme la
graine ne dépend que de l'identifiant de l'instrument, un script portant le même identifiant qu'un instrument de la bibliothèque voit les mêmes trajectoires.

\omcode{code/firm/payoffdsl/firm_payoffdsl.py}{403}{429}{Le moteur de script~: échéancier, trajectoires de la bibliothèque, actualisation sur la
courbe de l'instrument, et l'évaluateur.}\label{lst:pl:pricing-library-architecture:engine}

\begin{table}[ht]
\centering\small
\begin{tabular}{@{}>{\raggedright\arraybackslash}p{2.8cm}>{\centering\arraybackslash}p{1.8cm}>{\centering\arraybackslash}p{1.8cm}>{\centering\arraybackslash}p{1.8cm}>{\centering\arraybackslash}p{1.8cm}>{\raggedright\arraybackslash}p{2.5cm}@{}}
\toprule
produit & script & erreur type & biblio\-thèque & écart (e.t.) & moteur de la bibliothèque \\
\midrule
call européen & 11.3953 & 0.0570 & 11.3485 & 0.82 & analytique \\
call asiatique (12 fixings mensuels) & 6.8012 & 0.0327 & 6.8012 & 0.00 & Monte Carlo \\
call down-and-out quotidien & 11.0373 & 0.0568 & 11.0226 & 0.26 & forme fermée (BGK) \\
autocallable worst-of & 95.8750 & 0.0717 & 95.8750 & 0.00 & Monte Carlo \\
autocallable, knock-in lookback & 94.2474 & 0.0726 & --- & --- & aucun \\
\bottomrule
\end{tabular}
\caption{Scripts confrontés aux moteurs propres de la bibliothèque, 100\,000 trajectoires. Face à un moteur analytique ou à forme fermée, l'écart est une erreur Monte Carlo
inférieure à une erreur type~; face au moteur Monte Carlo de la bibliothèque, sur nombres aléatoires communs, elle est nulle, car les trajectoires et le
calcul du payoff sont les mêmes.}\label{tab:pl:pricing-library-architecture:validation}
\end{table}

\begin{example}[Prix dès lundi]\label{ex:pl:pricing-library-architecture:monday}
L'autocallable scripté vaut 95.8750 pour 100, et le moteur autocallable propre à la bibliothèque donne 95.8750~: l'écart vaut 0.00
erreur type, car les deux évaluent le même payoff sur les mêmes trajectoires. Leurs deltas par rapport au premier sous-jacent sont égaux aussi, 0.1518 par unité de spot
pour 100 de notionnel, puisque les grecques bump-and-reprice de la bibliothèque fonctionnent sur tout moteur. La variante lookback du client, que la bibliothèque ne peut pas exprimer,
vaut 94.2474 ($\pm0.0726$) avec la même mécanique. L'évaluateur compilé rend le moteur de script aussi rapide qu'un moteur à tableaux
écrit à la main~; l'interpréteur, 22 à 37 fois plus lent, reste la référence qui le vérifie.
\end{example}

\section{Recalcul paresseux et observateurs}

\begin{definition}[Motif observateur]\label{def:pl:pricing-library-architecture:observer}
Le \emph{motif observateur}\index{motif observateur} permet à un objet (l'observable) de tenir une liste de dépendants (les observateurs) et de les notifier
lorsqu'il change, sans savoir qui ils sont~; une courbe observe ses cotations, un prix observe sa courbe.
\end{definition}

\begin{definition}[Recalcul paresseux]\label{def:pl:pricing-library-architecture:lazy}
Le \emph{recalcul paresseux}\index{recalcul paresseux} marque un résultat comme périmé à l'arrivée d'une notification et ne le recalcule que lorsque le résultat
est demandé à nouveau, de sorte qu'une rafale de changements de marché ne coûte qu'un recalcul par lecteur, au lieu d'un par changement.
\end{definition}

La \texttt{LazyObject} de QuantLib est l'implémentation la plus connue~: un «~cadre de calcul à la demande et de mise en cache des résultats~» qui est à la fois
observateur et observable. Le \texttt{LazyPrice} du chapitre reprend la même idée en vingt lignes
(\cref{lst:pl:pricing-library-architecture:lazy})~: un call européen observant une cotation de taux est calculé une fois (11.35), la cotation change trois
fois, et la lecture suivante recalcule une fois (11.59) --- deux calculs pour quatre états du marché. La conception par instantanés de
\texttt{firm.pricing} est l'autre réponse au même problème~: rien n'est modifiable, donc rien ne devient périmé~; un nouvel instantané est une nouvelle question.
Les bibliothèques qui servent des outils interactifs privilégient les observateurs~; celles qui servent le calcul de risque par lots privilégient les instantanés.

\omcode{code/firm/payoffdsl/firm_payoffdsl.py}{469}{490}{Un prix paresseux~: périmé à la notification, recalculé à la demande, et lui-même observable
par ce qui en dépend.}\label{lst:pl:pricing-library-architecture:lazy}

\section{Liaisons vers d'autres langages}

Une bibliothèque écrite une fois est appelée depuis de nombreux endroits~: la recherche en Python, un service de pricing en C++, un tableur sur le desk d'un trader. Le chapitre~9
a mesuré ce que coûte une liaison par appel~; pour une bibliothèque, la leçon est de lier au niveau des requêtes entières --- un lot de transactions, un instantané
--- jamais au niveau d'une opération. La bibliothèque du livre~5 dispose de jumeaux C++20 et Rust de son cœur, testés face à la référence Python sur les
mêmes entrées~; un \omterm{def:pl:pricing-library-architecture:script}{script de payoff} est portable par construction, puisqu'il s'agit d'un texte que toute implémentation du langage peut pricer, ce qui explique pourquoi
les scripts circulent plus facilement que les classes entre les outils de front-office, les systèmes de risque et le modèle canonique de transaction du chapitre~21.

\section{Tutoriel~: un langage, un moteur, cinq produits}\label{tut:pl:pricing-library-architecture:tut}

\begin{tutorial}
\textbf{Objectif.} Écrire les produits sous forme de scripts, les pricer par la bibliothèque, les confronter aux moteurs de la bibliothèque, et mesurer les deux
évaluateurs. \textbf{État final~:} le \cref{tab:pl:pricing-library-architecture:validation} et la \cref{fig:pl:pricing-library-architecture:speed}.
\begin{tutsteps}
\item \textbf{Analyse}~: \texttt{firm\_payoffdsl.parse} sur chaque script~; en casser un et lire la position de l'erreur.
\item \textbf{Pricer}~: \texttt{firm.pricing.price} sur un \texttt{ScriptedInstrument}~; \texttt{greeks} dessus.
\item \textbf{Valider}~: \texttt{pl\_payoffdsl.validation()}.
\item \textbf{Mesurer}~: \texttt{bench\_payoffdsl.py}.
\item \textbf{Observer}~: \texttt{lazy\_demo()}.
\end{tutsteps}
\textbf{Pour aller plus loin.} Ajouter une fonction \texttt{worst} sur des sous-jacents nommés et un lookback quotidien entre les dates d'observation~; compiler un
script en source Python plutôt qu'en fermetures, et mesurer de nouveau.
\end{tutorial}

\section{Construction~: le langage de script de payoff}\label{bld:pl:pricing-library-architecture:payoffdsl}

\begin{build}
\bfield{Rôle} Des produits pour lesquels la bibliothèque n'a pas de classe, pricés, bumpés et calculés par lots comme ceux qu'elle possède, à partir de scripts qui peuvent être stockés,
relus et transférés d'un système à l'autre.
\bfield{Interface} \texttt{parse}, \texttt{schedule}, \texttt{Interpreter}, \texttt{compile\_program}, \texttt{ScriptedInstrument},
\texttt{ScriptEngine} (enregistré pour le modèle de Black--Scholes), \texttt{ScriptError}~; \texttt{Quote}, \texttt{LazyPrice}.
\bfield{Règles} Les trajectoires et graines de la bibliothèque, jamais celles du moteur~; actualisation sur la courbe de l'instrument~; l'interpréteur et
l'évaluateur compilé s'accordent sur chaque trajectoire~; les erreurs désignent une ligne et une colonne.
\bfield{Tests d'acceptation} \texttt{code/firm/payoffdsl/tests/}~: analyse et positions d'erreur~; l'ordre de l'échéancier des événements~; l'interpréteur égal
à l'évaluateur compilé sur un script dépendant du chemin~; un européen scripté égal au prix Monte Carlo de la bibliothèque et proche de son
prix analytique, avec son delta~; les prix paresseux et les observateurs enchaînés.
\bfield{Pour aller plus loin} Plusieurs sous-jacents nommés~; grecques adjointes (AAD) à travers le script compilé (livre~4, chapitre~28)~; une bibliothèque de scripts avec
des produits versionnés et relus.
\end{build}

\begin{omsources}
\item Documentation de référence de QuantLib, \texttt{LazyObject}.
\item Livre~5, chapitres 4, 18 et 28 (grecques, autocallables, la bibliothèque)~; livre~4, chapitre~26 (Monte Carlo).
\end{omsources}

\section{Exercices}

\begin{exercise}[$\star$]\label{exo:pl:pricing-library-architecture:1}
Écrire une option digitale d'achat qui verse 10 à l'échéance si le spot est supérieur à 105, sous forme de \omterm{def:pl:pricing-library-architecture:script}{script de payoff}.
\end{exercise}

\begin{exercise}[$\star$]\label{exo:pl:pricing-library-architecture:2}
Pourquoi le prix de l'autocallable scripté est-il exactement celui de la bibliothèque, alors que l'européen scripté s'écarte du prix analytique de 0.82
erreur type~?
\end{exercise}

\begin{exercise}[$\star$]\label{exo:pl:pricing-library-architecture:3}
Dans le script de l'autocallable, pourquoi le bloc du coupon doit-il précéder le bloc de l'autocall à la même date~?
\end{exercise}

\begin{exercise}[$\star\star$]\label{exo:pl:pricing-library-architecture:4}
Comment l'évaluateur compilé traite-t-il \texttt{stop}, et que se passerait-il s'il se contentait de sauter le reste du bloc~?
\end{exercise}

\begin{exercise}[$\star\star$]\label{exo:pl:pricing-library-architecture:5}
Pourquoi l'accélération de l'évaluateur compilé croît-elle avec le nombre de trajectoires, puis se stabilise-t-elle~?
\end{exercise}

\begin{exercise}[$\star\star$]\label{exo:pl:pricing-library-architecture:6}
Observateurs ou instantanés~: lequel choisiriez-vous pour l'écran de pricing d'un trader, et lequel pour le calcul de risque de nuit~?
\end{exercise}

\begin{exercise}[$\star\star\star$]\label{exo:pl:pricing-library-architecture:7}
\emph{Programmation.} Pricer l'autocallable lookback avec un niveau de knock-in à 50~\% puis à 70~\%. Comment le prix évolue-t-il, et que dit cela
de la sensibilité du produit à la barrière~?
\end{exercise}

\begin{exercise}[$\star\star\star$]\label{exo:pl:pricing-library-architecture:8}
\emph{Trouver la faille.} «~Notre moteur de script tire ses propres nombres aléatoires, donc les prix des scripts sont indépendants de la bibliothèque et
constituent un bon contrôle croisé.~»
\end{exercise}

\section{Problème~: prêt pour lundi}

\begin{problem}[{Problème du week-end --- un produit pour lequel la bibliothèque n'a pas de classe}]\label{pb:pl:pricing-library-architecture:1}
\texttt{firm.pricing}, \texttt{firm.payoffdsl} et les cinq produits du chapitre.

\medskip\noindent\textbf{Partie I --- La bibliothèque.}
\begin{enumerate}
\item Quelles sont les quatre couches de la bibliothèque, et qu'est-ce qui les relie~?
\item Que répond un \omterm{def:pl:pricing-library-architecture:market}{objet de marché}, et pourquoi un instantané immuable~?
\item Pourquoi les prix après choc ne diffèrent-ils que par le choc lui-même~?
\item Que doit fournir un moteur pour être enregistré~?
\item Pourquoi un système de risque n'a-t-il besoin que de la fonction de repricing~?
\end{enumerate}

\medskip\noindent\textbf{Partie II --- Le langage.}
\begin{enumerate}[resume]
\item Que peut dire un script, et que voit chaque bloc~?
\item Comment l'échéancier des événements est-il extrait~?
\item En quoi l'interpréteur et l'évaluateur compilé diffèrent-ils~?
\item En quoi le produit du client diffère-t-il de l'autocallable simple dans le script~?
\item Que se passe-t-il pour une erreur de syntaxe~?
\end{enumerate}

\medskip\noindent\textbf{Partie III --- Les chiffres.}
\begin{enumerate}[resume]
\item Donner les prix du script et de la bibliothèque pour les cinq produits, ainsi que leurs écarts en erreurs types.
\item Pourquoi deux écarts sont-ils exactement nuls~?
\item Quels sont les deltas de l'autocallable par le script et par la bibliothèque~?
\item Quelle est l'accélération de l'évaluateur compilé~?
\item Que compte la démonstration du prix paresseux~?
\end{enumerate}

\medskip\noindent\textbf{Partie IV --- Le verdict.}
\begin{enumerate}[resume]
\item Énoncer le \emph{résultat nommé}~: le prix et le delta de l'autocallable scripté face au moteur d'autocallable de la bibliothèque sur
nombres aléatoires communs, leur écart en erreurs types, et l'accélération du script compilé par rapport au script interprété.
\item Que faudrait-il exiger avant qu'un produit scripté soit négocié~?
\item Quand un produit doit-il avoir sa propre classe et son propre moteur plutôt qu'un script~?
\item Comment faudrait-il stocker et versionner les scripts~?
\item En une phrase~: qu'apporte un langage de payoff à une bibliothèque de pricing~?
\end{enumerate}
\end{problem}

\section{Questions d'entretien}

\begin{interviewq}[$\star$ \iqroles{développeur}]\label{iq:pl:pricing-library-architecture:1}
Pourquoi séparer les instruments, les modèles et les moteurs dans une bibliothèque de pricing~?
\end{interviewq}

\begin{interviewq}[$\star\star$ \iqroles{développeur, chercheur}]\label{iq:pl:pricing-library-architecture:2}
Comment pricerez-vous un nouvel exotique d'ici lundi sans écrire de nouveau moteur~?
\end{interviewq}

\begin{interviewq}[$\star\star$ \iqroles{développeur}]\label{iq:pl:pricing-library-architecture:3}
Comment rendez-vous un évaluateur de payoff Monte Carlo rapide en Python~?
\end{interviewq}

\begin{interviewq}[$\star\star$ \iqroles{développeur}]\label{iq:pl:pricing-library-architecture:4}
Qu'est-ce que l'\omterm{def:pl:dataframes:eager}{évaluation paresseuse} dans une bibliothèque de pricing, et quand est-ce une mauvaise idée~?
\end{interviewq}

\begin{interviewq}[$\star\star\star$ \iqroles{développeur}]\label{iq:pl:pricing-library-architecture:5}
Concevez un \omterm{def:pl:pricing-library-architecture:script}{langage de script de payoff} pour un desk de produits dérivés.
\end{interviewq}
